\section{Introduction}\label{sec:introduction}
Large language model (LLM) applications such as GPT~\cite{chatgpt}, Gemini~\cite{gemini}, and Claude~\cite{claude} increasingly incorporate web search functions, providing internet users with new ways of accessing information~\cite{webgpt-arxiv21, geo-apmr+24, zhai2024llm-and-future-of-information-retrieval}. 
Information retrieval has become a primary use case for LLMs, with a survey in September \(2025\) showing that about \(30\%\) of all prompts involve information seeking~\cite{how-people-use-chatgpt-nber25}. 
Systems that perform web searches and generate answers are known as \emph{generative engines} (GEs)~\cite{geo-apmr+24}. 
Web search services are rapidly transitioning from traditional search engines to GE-based systems~\cite{openai-chatgpt-search24, google-gemini-search24, perplexity-comet25, dia-browser25, microsoft-bing-chat23, anthropic-multiagent25}, changing the way users access and interpret web contents.

GEs not only alter the delivery of information but also reshape its nature. 
Whereas conventional search engines direct users to primary information sources such as official documents, GEs provide synthesized answers that summarize and interpret web content. 
LLMs also exhibit bias~\cite{bias-and-fairness-in-large-language-models-cl24, on-the-dangers-of-stochastic-parrots-facc21, holisticevaluationlanguagemodels-arxiv22, bbc-ai-search-engines-2024}, meaning that users receive secondary or higher-order information produced by LLMs rather than directly reading primary sources, similar to human-to-human information delivery~\cite{the-spread-of-true-and-false-news-online-science2018}.
This transformation enhances accessibility but centralizes the selection and interpretation of web content in the GE; moreover, major GEs are black-box systems, introducing vulnerabilities such as inaccuracy of answers.

GEs cite content published on the web by various publishers, including attackers with malicious intent or for profit; therefore, they remain vulnerable to poisoning attacks that inject malicious content into the web and produce attacker-intended answers~\cite{poisonedrag-arxiv24, poisoning-web-scale-training-datasets-is-practical-sp24}. 
When GEs generate answers lacking factuality, the cause is not only LLM hallucinations~\cite{sirens-song-in-the-ai-ocean-coli25, survey-of-hallucination-in-natural-language-generation-acmcsurv23} but also the lack of factuality in the web content obtained through web searches~\cite{retrieval-helps-or-hurts-deeper-arxiv24, poisonedrag-arxiv24, a-comprehensive-survey-on-trustworthy-recommender-systems-arxiv22}.
Web content can be published and modified by any user with any intent~\cite{poisoning-web-scale-training-datasets-is-practical-sp24, poisonedrag-arxiv24}.
Prior studies formalize poisoning attacks that place disinformation and misinformation on the web to generate attacker-intended answers on retrieval from external databases and answer generation systems including GEs because these systems can cite such contents and generate text that reflects them~\cite{poisoning-web-scale-training-datasets-is-practical-sp24, poisonedrag-arxiv24, yu2025safetydevolutionaiagents}.
Specific to GEs, \emph{Generative Engine Optimization} (GEO) presents patterns of which styles of web content GEs prefer to cite in answers~\cite{geo-apmr+24}.
Although GEO is beneficial for primary information providers that seek to increase exposure of their web content as citations, it can also facilitate poisoning attacks.

Existing evaluation criteria aim to assess how faithfully the cited content is reflected in the answers generated by the retrieval and generation systems ~\cite{citeeval-acl25, ragas-eacl24, ares-naacl24, evaluating-verifiability-in-generative-search-engines-emnlp23, enabling-llms-to-generate-text-with-citations-emnlp23, ranking-generated-summaries-by-correctness-acl2019, evaluating-the-factual-consistency-of-abstractive-text-summarization-emnlp20, alignscore-acl23, feqa-acl20}.
These studies analyze textual and semantic consistency between the cited web content and generated answers.
However, these studies alone cannot capture the attack surface of GEs to PoisonedRAG attacks.
% These studies aim to verify the consistency between the content cited by GEs and the answers generated by GEs, but they do not aim to determine whether the answers may be contaminated by malicious content.
This is because these studies targeted systems called \emph{RAG systems}, which retrieve from curated, fixed, and closed external databases to generate answers.
In contrast, GEs treat the web as an external database, which is open to anyone and has dynamism, allowing attackers to easily inject malicious content that GEs can cite.
Due to these differences in characteristics, evaluation criteria proposed for \emph{RAG systems} cannot adequately assess the reliability and safety of GEs.

To address this gap, we introduce novel evaluation criteria that focus on the publisher attributes of citation in answers. 
Our goal is to reveal the bias of GEs in which publisher attributes GEs preferentially cite and how faithfully they reflect the content in generated answers.
First, we classify the publisher attributes of citations in answers following the steps: (1) classification of publishers into primary and secondary information sources, (2) classification of secondary information sources into fine-grained publisher categories using a \emph{LLM-as-a-Judge} method~\cite{survey-llm-as-a-judge-arxiv24, judging-llm-as-a-judge-with-mt-bench-and-chatbot-arena-arxiv23} based on URL and WHOIS data, (3) we classify the category of publishers into \emph{content-injection barriers}, representing the practical difficulty of publishing content with specific publisher authority, specifically into low-, medium-, or high-barrier levels to reveal the attack surface of GEs and the difficulty for poisoning attacks to succeed.
Second, we propose two evaluation criteria: (a) distribution of content-injection barriers in answers and (b) distribution of semantic consistency by content-injection barriers.
Using our criteria, we reveal the bias of GEs in their citation selection and content reflection in answers from the perspective of the authority of publishers, and characterize the attack surface of GEs and the difficulty for poisoning attacks to succeed.

To reveal the threat of poisoning attacks on representative GEs (OpenAI GPT-5, Claude 4 Sonnet, and Gemini Flash 2.0 with search mode enabled), we apply these evaluation criteria to information-retrieval questions in the political domain across Japan and the United States (U.S.), using \(280\) questions and \(4{,}200\) generated answers.
Our results show that citations to official party websites (primary information) constitute approximately \(60\%\)--\(65\%\) of citations in Japan but only \(25\%\)--\(45\%\) in the U.S. and low-barrier sources that can be published with only registration (Reddit, X, personal blog, etc.) account for approximately \(30\%\) of citations in answers.
We further find that citations from low-barrier sources tend to have lower semantic consistency with the actual answer content than medium- and high-barrier sources, yet they still influence GE answers despite their lower semantic consistency.

In summary, our study makes the following contributions:
First, we propose new evaluation criteria with publisher attributes of citations in generated answers by GEs beyond existing evaluation criteria for \emph{RAG systems}.
Second, we find that GEs cite primary sources in \(60\%\)--\(65\%\) of Japanese political answers and \(25\%\)--\(45\%\) of U.S. ones, while low-barrier sources account for approximately \(30\%\) of citations, characterizing the attack surface of GEs and the threat of PoisonedRAG attacks, as GEs tend to cite low-barrier sources.
